\begin{pastexam}[label={px:ust-113-ee-ds-05}]{臺聯大電機類 113 學年度 第 5 題}
（單選；每題 3 分，答錯不倒扣）

What should be added to the blank so that the following function can correctly reverse a singly
linked list.
\begin{minted}[fontsize=\small, style=perldoc, breaklines]{c}
struct node {
    int data;
    struct node * next;
};
static void reverse(struct node ** head) {
    struct node * prev = NULL;
    struct node * current = *head;
    struct node * next;
    while (current != NULL) {
        next = current->next;
        current->next = prev;
        prev = current;
        current = next;
    }
    ________________________
}
\end{minted}

(A)~\texttt{*head = prev;} \qquad (B)~\texttt{*head = current;} \qquad (C)~\texttt{*head = next;}\\
(D)~\texttt{*head = NULL;} \qquad (E)~None of the above
\end{pastexam}

\begin{pastsoln}
答案是 (A)。

這是資結書虛擬碼~5.12 的就地反轉（資結書第~5.5~節）。 迴圈每一輪把 $\mathit{current}$ 的指標反過來指向
$\mathit{prev}$，然後兩個指標一起往前一格。 不變量是：$\mathit{prev}$ 是「已經反轉好的那一段」的開頭，
$\mathit{current}$ 是「還沒處理的那一段」的開頭。 迴圈結束時 $\mathit{current} = \texttt{NULL}$，整條都反轉
好了，新的開頭是 $\mathit{prev}$（原本的最後一個節點），所以要寫 \texttt{*head = prev;}。

例：$1 \to 2 \to 3$。
\begin{center}
\begin{tabular}{llll}
\toprule
輪 & 反轉好的部分（從 $\mathit{prev}$ 起） & $\mathit{current}$ & $\mathit{next}$ \\
\midrule
開始 & （空） & 1 & \\
1 & $1$ & 2 & 2 \\
2 & $2 \to 1$ & 3 & 3 \\
3 & $3 \to 2 \to 1$ & \texttt{NULL} & \texttt{NULL} \\
\bottomrule
\end{tabular}
\end{center}

逐項：
\begin{itemize}
\item (A) 對。
\item (B) 迴圈結束時 $\mathit{current}$ 一定是 \texttt{NULL}，等於把串列弄成空的，節點全部遺失。 錯。
\item (C) 最後一輪 $\mathit{next} = \mathit{current}{\to}\mathit{next} = \texttt{NULL}$，同 (B)；串列本來
就是空的時候 $\mathit{next}$ 根本沒有設定過，是未初始化的值。 錯。
\item (D) 同樣把串列變成空的。 錯。
\item (E) (A) 是對的。 錯。
\end{itemize}
\end{pastsoln}
